\capitulo{6}{When Is Getting It Right Enough?}
I have two groups of learners at the workshop, grading tickets. Whenever the two faces match, everyone writes A. One group thinks nine out of ten such tickets come from A. The other thinks six out of ten do. If the test asks only for the letter on each ticket, I can swap the two groups' answer sheets and nobody will notice.

You know a way to tell them apart. Walk back into the room with the ticket that costs R\$\,8.80 and pays ten reais if the machine was A. One group expects to collect nine reais and buys it. The other expects six and passes. Suppose, though, that all the workshop wants is to route tickets into two drawers, and that a mistake costs the same in either direction. For that job, how much is it worth to distinguish two bets that open the same gate?

\secao{What Do You Charge for a Right Answer?}
Let's keep the faces, the labels, and even the form of the function we use to classify them. As in the last chapter, $z=+1$ means the faces match, and $z=-1$ means they differ. Our score will be $f_w(z)=wz$. Give A the label $y=+1$ and B the label $y=-1$. The symmetry of our workshop lets us use the same coefficient for both kinds of ticket. On matching faces, $w=0.4$ gives A, and so does $w=4$. What would make one number earn its place over the other?

Logistic regression has an answer. Its score measures the log odds of A. By the time we grade the test, though, we know which letter came up. A score of $-2$ when the letter is B deserves as much credit as $+2$ when the letter is A. We can turn the ruler around for the B tickets. Multiplying the score by the label $y$ leaves the reading alone for A and reverses its sign for B. Call this reading in favor of the letter that actually came up $m=yf_w(z)$. For both right answers in the example, $m=2$.

What probability did the learner give the right letter? On the ruler we just turned around, the odds in its favor are $e^m$ to one. Of the $1+e^m$ parts, $e^m$ belong to that letter. Its share is therefore $p_{\rm right}=1/(1+e^{-m})$, as in the previous chapter. To work out the charge, look at the denominator. It tells you what you would have to multiply that share by to get back to one. The logarithm of that factor is the price we set in Chapter 2, this time in natural logarithm units,
\[
             \ell_{\rm log}(m)=-\ln p_{\rm right}=\ln(1+e^{-m}).
\]
Even a learner who got the letter right still owes something. Among the tickets with matching faces, the bill is smallest when we quote the fraction of A tickets in the pile, as we saw in Chapter 4. Nine A tickets and one B ticket call for odds of nine to one, or $w=\ln 9$. If the pile has six A tickets and four B tickets, the odds are three to two and the score drops to $\ln(3/2)$. With no penalty, the fractions for the tickets with different faces are mirror images of these, so the same coefficient works for both groups.

Let me propose another way to grade the test. I'll charge one point for a tie, $m=0$. For a step in the right direction, I'll take a point off the bill. For a step in the wrong direction, I'll add a point. The bill has a floor. Once you reach $m=1$, I charge nothing more on that ticket. A right answer with $m=0.6$ therefore costs $0.4$ points and a wrong one with $m=-0.6$ costs $1.6$. This rule is called the \emph{hinge loss},
\[
                  \ell_{\rm h}(m)=\max(0,1-m).
\]
We call the product $m$ the \emph{functional margin}. Let's see what the new bill does to those nine right answers.

\perdasmargem

Put ten tickets with matching faces on the table, nine from A and one from B. At $w=0$ the hinge bill comes to ten points. Move from zero to one. You take a point off nine tickets and pay an extra point on one. You still owe two points. Now suppose you want to go from one to two. Which of the nine right answers will still earn you a discount? None of them. The only bill that changes belongs to the wrong answer, and it goes up by one point. The loss is linear on each stretch, and moving to the negative side would favor the minority. The best place to stop is $w=1$.

Trade three A tickets for three B tickets. You now have six against four. Up to $w=1$, the six discounts still cover the four extra charges. Beyond one, all that happens is that the four mistakes get more expensive. The best place is still $w=1$. I could hand you the winning coefficient, explain exactly how I chose it, and ask what I ought to pay for the next ticket. Could you tell whether I had the pile with nine A tickets or the one with six? The gate reads the positive sign and gives the same answer for both piles. Now suppose someone feeds the score of one through the logistic sigmoid and announces a probability of about $73\%$ for A. That gives us a price of R\$\,7.31 for both piles. One pile suggests nine reais and the other suggests six. I could defend R\$\,7.31 if I found that frequency in fresh tickets with the same score. But that promise took an extra assumption about what the number one meant. The hinge loss only charged us for the margin.

\secao{Where the Margin Comes From}
To get to a \emph{support vector machine}, or SVM, we have one more choice to make. The workshop inspector brings us a caliper. The gate now sorts parts by their measured diameter. In the batch we fit to, the largest B part measures $4.8$ mm and the smallest A part measures $5.2$ mm. Anything above the threshold goes to A. I set my threshold at $5.15$ mm. You set yours at $5.00$ mm. We test the parts one by one, and we both get every one right. Before placing an order for your gate or mine, what would you like to ask the inspector?

He takes the part that measured $5.2$, lifts it out of the caliper, and clamps it again. Just changing its position, he tells us, can make two readings differ by as much as $0.10$ mm. Now there is a difference between our two perfect scores. A second reading of $5.10$ sends the part through my gate to B. At your gate it still has room to spare. The same reasoning works from the other side. Among the boundaries that got the whole batch right, the one at $5.00$ mm gives the nearest part on either side the most clearance. The person who knew how much that clearance was worth was holding the caliper.

Why didn't we simply reward a high score on the right side? If our ruler is $f(\boldsymbol u)=\boldsymbol w^\top\boldsymbol u+b$, I can double $\boldsymbol w$ and $b$ without moving the boundary $f=0$ by so much as a millimeter. The scores have doubled. The part is just as far from the boundary as before. To connect the two scales, imagine that walking one millimeter perpendicular to the boundary adds four points on the display. To pick up one point, you need to move a quarter of a millimeter. The norm $\lVert\boldsymbol w\rVert$ plays the part of those four points per millimeter. Take the displacement along each axis, multiply it by its weight, and add the results. This measures how far you went in the direction of $\boldsymbol w$, scaled by its length. In that direction, perpendicular to the boundary, a distance $d$ adds $d\lVert\boldsymbol w\rVert$ to the score. The line where the score is $+1$ therefore lies $1/\lVert\boldsymbol w\rVert$ from the boundary. The line where it is $-1$ lies just as far away on the other side.

When we can give every example $y_i f(\boldsymbol u_i)\geq1$, choosing the smallest norm $\lVert\boldsymbol w\rVert$ gives us the widest strip. With our caliper, that puts the threshold halfway between the two nearest parts. But for a ticket with two heads, we know the very same input can carry either label, A or B. No straight line can put them on different sides. Let's strike a bargain. We want to widen the strip and we're willing to pay for the parts that fall inside it or on the wrong side. If I double the number of parts by copying each one, I want the average cost of the errors to stay the same. That's why I divide the bill for the parts by $n$. On the other pan of the balance I choose to charge for the square of the norm. Doubling the scale makes that charge four times as large. I have made a choice about what the width of the strip is worth. With both charges included, the cost minimized by a linear \emph{soft-margin} SVM is
\[
 \frac1n\sum_{i=1}^n \max\{0,1-y_i f(\boldsymbol u_i)\}
       +\frac{\lambda}{2}\lVert\boldsymbol w\rVert^2,\qquad \lambda>0.
\]
The first term is the charge we applied to our two piles. The second charges for scaling up the coefficients even when we leave the boundary exactly where it was. A larger $\lambda$ puts more weight on the norm. The location of the boundary depends on the parts that reach the margin or violate it. Among them are the \emph{support vectors}. The value $w=1$ in our envelope example came from the loss alone. In fitting an SVM, the penalty has a say too.\footnote{Corinna Cortes and Vladimir Vapnik develop the classifier with a margin and margin violations in \href{https://doi.org/10.1007/BF00994018}{\emph{Support-Vector Networks}} (1995). For the formulation in terms of a loss and a penalty, see T. Hastie, S. Rosset, R. Tibshirani, and J. Zhu, \href{https://jmlr.org/papers/v5/hastie04a.html}{\emph{The Entire Regularization Path for the Support Vector Machine}} (2004).}

At the workshop I would test the SVM if I wanted the sorting decision to hold up under repeated measurements and could justify the ruler I used to measure distance. If I give the machine both diameter and weight, changing grams to milligrams already changes the geometry of the raw data. Nobody learned anything new about the parts by changing units. It falls to me to decide what a step along each axis is worth. To check my choice, I would set aside parts of known origin and measure them again. The clearance I found on yesterday's parts interests me only if it helps route the ones coming through today.

Now change the job. We want to decide whether a part is worth inspecting. An inspection costs two reais and prevents a loss of ten reais if the part came from B. Call the risk of B at that measurement $r$. An inspection is worth doing when $10r>2$. At an inspection price of four reais, the cutoff moves from $20\%$ to $40\%$. A part with a $30\%$ risk changes destination without anybody touching the caliper. Here I would start with logistic regression. I can take its risk estimate to the person who knows the prices. There is one detail. In the batch where A and B separated perfectly, who measured the risk for a $5.05$ mm part? No part had that diameter. A steeper and steeper logistic curve can keep improving the fit to the letters in the batch. Before spending money on the strength of its number, I would find labeled parts near the decision threshold and check whether the announced risks match what happens. We can also learn to calibrate an SVM score. For that I would set aside labeled data from the population where I plan to use the score and measure the frequency of B there.

\secao{Can We Learn When the Answer Stays the Same?}
Set both programs to $w=2$. Give each of them a ticket with matching faces that turns out to be an A. Both got it right. Now slide your finger to the right on the loss graph, from $2$ to $2.01$. On the hinge curve, your finger traveled along the flat stretch. The logistic curve still slopes down. Just by looking, you can tell which learner has a reason to increase $w$ after seeing this ticket.

How much does it still stand to gain? At $m=2$, the logistic charge is about $0.1269$. At $2.01$ it is $0.1257$. A step of one hundredth saved about $0.0012$ points. The local slope must be close to $-0.12$ points per unit of margin. We can understand the size of that number by looking at the bet. In $\ln(1+e^{-m})$, the one stays put. If $m$ increases by one hundredth, the term $e^{-m}$ gets multiplied by $e^{-0.01}$, about $0.99$. It loses nearly a hundredth of itself. As we take smaller and smaller steps, the amount it loses per unit step approaches $e^{-m}$ itself. The logarithm measures that change relative to the total $1+e^{-m}$. The slope is therefore
\[
 \ell'_{\rm log}(m)=-\frac{e^{-m}}{1+e^{-m}}=-(1-p_{\rm right}).
\]
The learner who put close to $88\%$ on the right letter still had nearly $12\%$ on the wrong side. That leftover share is what the ticket contributes to the adjustment. A \emph{gradient step} tells us to move against the slope to bring the bill down. If $\eta>0$ controls the step size, then, since $m=w$ here, the ticket adds $\eta/(1+e^2)$ to $w$. On the hinge curve, a slope of zero adds zero. Suppose we also charge $\lambda w^2/2$ for the coefficient. Why does a large $w$ feel a stronger pull back toward zero? Picture $w^2$ as the area of a square with side $w>0$. Extend the side a little and you add two strips of length $w$ and a tiny corner that shrinks with the square of the step. The area grows at a rate of $2w$. Multiply by the rate we charge, $\lambda/2$, and you get a slope of $\lambda w$. That part subtracts $\eta\lambda w$ from both updates. We can now tell how much of the turn of the dial came from the new letter and how much came from a charge that was already there.

We could have made a different choice at the outset. For each kind of ticket, keep just the counts of A and B. When a new letter arrives, add one to its count. With the same notebook I can tell you which letter occurs more often, calculate the fraction of A tickets, or choose the score that minimizes either loss in that column. Changing the question I asked the notebook cost me a fresh calculation when I read it. The instruction to add one to the count still worked.

On the first page of this book, the state of the learner was a distribution, and the rule that updated it was called $\mathcal F$. The notebook has grown. At one point it held counts and at another it held the coefficients of a line. If $S_t$ is what we keep, $d_t$ the new datum, and $Q$ the question, we now have room for two operations,
\[
        S_t=U(S_{t-1};d_t),\qquad
        \text{answer}=R_Q(S_t).
\]
When the state is a distribution over hypotheses, $U$ can be our $\mathcal F$, and $R_Q$ can simply read off the probability of A. We needed the distinction once we could ask the same state different questions and give the same datum to rules that handle it differently.

\begin{quadro}{What I Chose Before Opening the Notebook}
\itemq{To place a bet} In logistic regression, I choose the features that make up the log odds and charge according to the probability given to the letter that was revealed. Before I use the result as a risk, I want to see how often errors occur among the parts whose fate depends on that risk.
\itemq{To separate} With an SVM, I choose a representation, a ruler for measuring distance, and a price for margin violations. A correctly classified example beyond the margin has already paid its share of the loss.
\end{quadro}

The pile with nine A tickets and one B ticket is still on the table. You know only those counts and you bet according to the observed frequency. I bring another ticket with matching faces. We turn it over. B. You cross out the ten at the bottom of the column and write eleven. I ask about the next letter. You say A, just as you did before. I hold out the ticket that pays ten reais and ask R\$\,8.80 for it. It is now worth $10\times9/11$ to you, about R\$\,8.18. You leave my hand hanging there.

If I had written down only the letter you predicted, I could have spent the whole afternoon without noticing that you had learned something.

\begin{quadro}{What to take away from this chapter}
\itemq{1} The loss we choose defines what counts as improvement, even when the predicted letter stays the same. Logistic loss and hinge loss charge for different things.
\itemq{2} In an SVM, the loss and the norm penalty work together. The geometric margin depends on the representation and the scale of the measurements.
\itemq{3} A classification score needs calibration checked against data to support decisions based on risk probabilities.
\itemq{What changed} We separated two operations: updating the learner's state and reading an answer from it. A class prediction can stay the same while the bet and the decision to buy change.
\itemq{To carry forward} If two notebooks announce $90\%$, should the next observation change both bets in the same way?
\end{quadro}
